\documentclass[11pt]{letter}
\usepackage[margin=1in]{geometry}
\usepackage[T1]{fontenc}
\setlength{\longindentation}{0pt}
\setlength{\indentedwidth}{\textwidth}
\signature{Tanishk Yadav\\Department of Finance and Risk Engineering\\New York University Tandon School of Engineering\\tanishkyadav@nyu.edu}
\address{}
\date{}
\begin{document}
\begin{letter}{The Editors\\Finance Research Letters}
\opening{Dear Editors,}

I submit ``Calm Months, Crash Months: Why Volatility Management Pays for Momentum and Not the Market'' for consideration in \textit{Finance Research Letters}.

Volatility-managed portfolios lose most of their appeal in real time, yet momentum keeps its gain. The letter explains why. Sorting stocks on past returns makes momentum's market beta a function of the market's own return over the formation window. Using Tweedie's formula, I derive this beta in closed form from two quantities measured in the cross-section of stocks, and test it on CRSP without fitting any parameter: the prediction tracks how momentum's beta moves (slope 0.92 from 1932 and 1.08 from 1963) and attains 79\% of the explanatory power of an in-sample regression. Months after market declines supply 71\% of the covariance through which volatility timing adds to momentum's return. If the market's premium rises with its variance, the same reasoning implies that the market should not be volatility-timed, and in the data it gains nothing from it.

The main test was specified in a plan committed before the CRSP data were downloaded. The manuscript has about 1,600 words of main text, three tables and three figures, and includes the required declaration on the use of generative AI. It is not under consideration elsewhere. An earlier version of the underlying project is posted on SSRN.

\closing{Sincerely,}
\end{letter}
\end{document}
